% !TeX root = ../main.tex
% Revised 2026-10-18: expanded primary-source citation traceback; implementation/results TODOs are intentional.
\section{Introduction}
\label{sec:intro}


Visual world models aim to predict how an environment evolves by generating future observations conditioned on current observations and actions~\cite{ha2018worldmodels,yang2024learninginteractiverealworldsimulators,hu2023gaia1generativeworldmodel}.
%
Building on advances in video generation~\cite{ho2022videodiffusionmodels,blattmann2023stablevideodiffusionscaling,yang2025cogvideoxtexttovideodiffusionmodels,kong2025hunyuanvideosystematicframeworklarge,wan2025wanopenadvancedlargescale}, recent models have enabled increasingly realistic and controllable visual rollouts driven by action inputs~\cite{bruce2024geniegenerativeinteractiveenvironments,valevski2025diffusionmodelsrealtimegame,he2026matrixgame20opensourcerealtime,zhang2026worldplay2extendingrealtimeinteractive,zhou2026astronexworld10realtimeinteractive} and camera trajectories~\cite{he2025cameractrlenablingcameracontrol,wang2024motionctrlunifiedflexiblemotion,xu2024camcocameracontrollable3dconsistentimagetovideo,zhang2026unifiedcamerapositionalencoding,huang2026solarwmopendatascalable,zhu2026sanawmefficientminutescaleworld}.
%
By predicting future observations under different actions or camera trajectories, these models provide a foundation for simulation, planning, and decision-making~\cite{yang2024learninginteractiverealworldsimulators,alonso2024diffusionworldmodelingvisual,gao2026dreamdojogeneralistrobotworld}.
%
Extending these capabilities to low-altitude unmanned aerial vehicles (UAVs) could aid in the training and evaluation of vision-based navigation systems by generating realistic onboard observations across diverse environments and flight trajectories~\cite{zhao2025airscapeaerialgenerativeworld,zhang2026aerialworldmodellonghorizon}.
%
However, realizing such low-altitude world models requires both flight data with accurate camera annotations and training mechanisms that align generated observations with prescribed camera trajectories.

These two requirements are closely linked, as training data with reliable camera annotations are a prerequisite for developing accurate camera-trajectory control mechanisms.
%
On the data side, public video datasets often lack consistent camera supervision due to missing trajectory annotations or discrepancies in coordinate conventions and translation scales~\cite{ju2024miradatalargescalevideodataset,zheng2025realcamvidhighresolutionvideodataset}.
%
Although SolarWM~\cite{huang2026solarwmopendatascalable} and SANA-WM~\cite{zhu2026sanawmefficientminutescaleworld} combine camera-motion estimation with scale alignment based on metric-depth predictions, their pipelines for video-only data still rely on predicted geometry rather than ground-truth reference for the recovered trajectory.
%
This reliance on predicted geometry can propagate pose and scale errors into the training annotations, hindering the model from learning accurate camera-trajectory control.
%
At the method level, camera-conditioned RGB generators use camera poses as conditioning signals for video synthesis~\cite{he2025cameractrlenablingcameracontrol,huang2026solarwmopendatascalable,zhu2026sanawmefficientminutescaleworld}.
%
However, such conditioning leaves scene geometry to be learned implicitly from RGB supervision, without explicit constraints or supervision on the underlying 3D structure.
%
These limitations motivate more reliable camera annotations and explicit supervision of geometry.

To this end, we propose \textbf{LAVi}, a framework for \textbf{L}ow-\textbf{A}ltitude \textbf{Vi}sual world modeling that couples ground-truth-guided camera annotation with joint RGB-depth generation.
%
On the data side, we assemble a corpus of 800,000 video clips spanning simulated environments, public video datasets, and real UAV flights.
%
It includes 470,000 clips collected in our simulation environment, with ground-truth camera trajectories and depth maps, together with self-collected UAV footage covering 300,000~m$^2$.
%
Beyond serving as training labels, these exact, metric-scale camera trajectories provide a reference for calibrating scale recovery and selecting pose estimators based on translational and rotational errors measured against ground truth.
%
The calibrated procedure and selected estimator are then integrated into a unified data construction pipeline that standardizes camera annotations across diverse video sources.
%
On the method side, we introduce explicit geometric supervision through joint RGB and depth generation under prescribed camera trajectories, encouraging geometrically coherent video generation under camera motion.
%
Supervising depth makes the arrangement and relative distances of visible surfaces explicit learning targets, providing structural cues for modeling parallax and visibility changes.
%
To realize this joint RGB-depth generation, we build upon Wan2.2-TI2V-5B~\cite{wan2025wanopenadvancedlargescale}, processing RGB and depth tokens through shared early Transformer blocks followed by modality-specific branches.
%
Training proceeds in two stages, beginning with camera-conditioned bidirectional learning using SCoPE~\cite{yin2026scope}, followed by a causal training stage that jointly aligns predictions with reference videos under ground-truth histories and matches self-generated rollouts to the reference distribution in feature space~\cite{xu2026directforcingsinglestagetraining}.
%
This enables few-step autoregressive generation along prescribed camera trajectories, providing a basis for a unified simulation environment with controllable flight scenarios, realistic visual observations, and explicit geometric information to aid the training and evaluation of UAV navigation systems.

% Keep this results paragraph and the contribution list compact under flush-bottom layout.
\begingroup
\setlength{\parskip}{0pt}
Experiments demonstrate that our camera annotation method achieves the lowest pose errors among the evaluated pipelines on data with ground-truth camera poses, while joint depth generation improves camera-control accuracy, enabling more faithful adherence to prescribed camera trajectories than existing camera-controlled world models.
Our main contributions are summarized as follows:
\begin{itemize}[topsep=2pt,partopsep=0pt,parsep=0pt,itemsep=0pt]
    \item We propose \textbf{LAVi}, a unified framework for low-altitude visual world modeling that integrates camera-annotated data construction with geometry-supervised video generation, enabling joint RGB-depth prediction along prescribed flight trajectories.
    \item We construct a corpus of 800,000 video clips spanning simulation, public datasets, and real UAV flights, including 470,000 simulated clips with exact metric camera trajectories and depth maps. These ground-truth trajectories guide scale calibration and pose-estimator selection for camera annotation and metric-scale alignment across sources.
    \item We develop a camera-control mechanism that combines trajectory conditioning with explicit geometric supervision through joint RGB-depth generation. This design improves camera-control accuracy, enabling more faithful adherence to prescribed camera trajectories than existing camera-controlled world models.
\end{itemize}
\endgroup
